\section{Introduction}\label{sec:intro}

The models that perform best on tabular data are tree ensembles whose reasoning is not
readable from their parameters~\cite{grinsztajn2022why}, and they are increasingly
deployed where decisions affect people. Regulation has followed: the EU AI Act obliges
providers of high-risk systems to make operation interpretable~\cite{euaiact2024}.
Two post-hoc methods dominate practice. \shap{} attributes a prediction to features
through cooperative-game theory~\cite{lundberg2017unified,lundberg2020local}; \lime{}
fits an interpretable surrogate in the neighbourhood of a single
prediction~\cite{ribeiro2016why}. Their foundations differ, one is deterministic on
trees while the other samples, and their failure modes are
unrelated~\cite{slack2020fooling,kumar2020problems}. A practitioner who runs both and
obtains two different top-five lists has no principled basis for choosing.

Prior comparisons do not settle this because they hold too much fixed: typically one
model family or one dataset, a locally defined fidelity metric, no hypothesis test, no
random-ranking baseline, and no statement of which output space the attributions occupy.
The gap is not a missing algorithm but a missing controlled comparison --- and that is
something an experiment can close.

This paper contributes:

\begin{enumerate}
  \item \textbf{A reproducible multi-axis benchmark.} Two explainers, three rank-fusion
    rules, and a random control across six classifiers and three public clinical
    datasets (\NumCells{} cells, \NumExplained{} instances), with eleven pre-registered
    hypotheses, non-parametric tests, Holm correction within family, effect sizes, and
    bootstrap intervals (Table~\ref{tab:hypotheses}).
  \item \textbf{A three-way decomposition of explanation stability} --- explainer
    randomness, input perturbation, and variance from retraining the model itself
    (Section~\ref{sec:rq2}). Reporting one stability number conflates these, and the
    distinction matters: an explanation can be bit-identical across runs and still
    fragile, because the model beneath it moves when the training sample does.
  \item \textbf{A hybrid consensus method, and an honest account of its limits.} We
    propose \emph{fidelity-weighted rank fusion}: the two explainers are combined in
    proportion to the faithfulness each demonstrates, with the weight fitted on a
    calibration split disjoint from the instances the method is later tested on, so it
    cannot win by construction. It is evaluated against two published fusion baselines ---
    mean-rank (Borda) and reciprocal-rank fusion~\cite{cormack2009reciprocal} --- and
    against a deliberately strict bar: beat the \emph{better} of \shap{} and \lime{} on
    the same cell, not their average, which anything would beat. Explanation aggregation
    itself is prior work~\cite{bhatt2020aggregating} and we claim no priority for the
    idea; what is new is the fidelity weighting, the leakage-safe calibration, and the
    controlled test. The outcome is negative, it was pre-registered as a valid result
    before running, and Section~\ref{sec:rq6} explains the mechanism.
\end{enumerate}

Section~\ref{sec:method} states the protocol, including four methodological traps that
silently invalidate comparisons of this kind. Section~\ref{sec:results} reports one
subsection per research question.
